\chapter{Interpretation}
This chapter provides a critical evaluation and discussion of the experimental results obtained from the implementation of the distributed robotic vision framework. By analyzing the system's performance holistically, this section identifies the underlying causes of the operational failures, highlights the theoretical and practical strengths of the architecture, outlines prominent technological limitations, and contrasts the proposed method against existing approaches in the field of mobile robotics.

\section{Discussion of Results}
Based on the empirical evidence gathered during the testing phases, it is evident that both implementation attempts faced significant challenges regarding real-time viability, as each approach was constrained by a distinct hardware or network bottleneck. The initial edge-computing attempt was halted by the raw processing limitations of the embedded hardware, while the subsequent network-distributed approach was heavily impacted by wireless data transmission latency.

The primary lesson derived from these results is that a successful person-following behavior requires a critical balance between high-level perception tasks and real-time execution loops. Even though the custom ROS 2 nodes for object detection and spatial calculation functioned as intended when evaluated independently, the systemic latency introduced by network distribution caused an immediate control loop delay. This latency manifested prominently during the searching phase, where the time-lagged coordinates led to an operational overshoot, causing the robot to rotate past the target before the feedback loop could stabilize. This outcome underscores that offloading computational workloads across distributed systems requires precise synchronization to prevent data delays from degrading active locomotive tracking.

\section{Strengths}
In a theoretically operational demonstration, the architectural design of this setup offers substantial advantages for mobile robot perception. The integration of a deep-learning-based computer vision model significantly elevates external environmental data collection by introducing semantic awareness. Unlike traditional distance rangefinders, this system possesses the ability to actively identify and classify specific objects, enabling the robot to adapt its behavioral states and safety boundaries contextually. Crucially, this vision pipeline is not intended to fully eliminate traditional sensors like LiDAR or mechanical bumpers; rather, it is designed to complement them, establishing a multi-modal sensor fusion array that increases overall system reliability. 

Furthermore, this setup provides a powerful foundation for space exploration and automated environment mapping. In scenarios where a novel facility needs to be mapped, a person-following framework is highly efficient, as a human operator can simply guide the robot through the building while simultaneous localization and mapping (\texttt{SLAM}) routines run concurrently in the background, minimizing the need for manual teleoperation.

\section{Limitations}
Despite these theoretical strengths, the technological and infrastructural limitations of the current implementation remain prominent. Although computer vision has experienced massive advancements in recent years, it remains heavily dependent on powerful or highly specialized processing hardware to run efficiently. This dependency was clearly demonstrated by the absolute failure of the Raspberry Pi~4 to execute even the smallest YOLO model at a usable frame rate, emphasizing the challenges of deploying modern AI models on standard embedded edge systems.

Additionally, the requirements of distributed network data distribution cannot be underestimated. Periodically transmitting uncompressed, high-frequency image messages and depth matrices in real time demands immense network bandwidth. Standard wireless Wi-Fi connections frequently struggle with or completely fail under this data load, introducing transport latencies that saturate the DDS middleware. Consequently, offloading vision tasks to an external workstation introduces a severe dependency on wireless network quality, compromising the real-time responsiveness of the robot's locomotion and causing the systemic overshoot observed during the searching phase.

\section{Comparison with Existing Approaches}
When contrasted against existing mobile robotics approaches, the proposed framework occupies a unique middle ground between centralized edge-computing platforms and cloud-based robotic setups. Traditional ROS-based person-following solutions often rely strictly on 2D LiDAR clustering, which is highly efficient and circumvents network bottlenecks but suffers from a complete lack of semantic intelligence and target cross-validation. 

Conversely, state-of-the-art commercial systems that implement advanced visual tracking typically bypass the network constraints observed in this thesis by utilizing dedicated local hardware acceleration—such as integrated onboard GPUs or utilizing the full processing potential of the camera’s internal VPU. By comparing these methodologies, it becomes clear that while the network-offloading approach via CycloneDDS developed in this research offers a highly cost-effective and modular solution for rapid software prototyping, physical industrial deployment ultimately requires localized hardware acceleration to eliminate the vulnerabilities associated with wireless data transmission and network-induced data latency.
